\section{Cyclic decomposition and ratios of determinants}
The cyclic projectors
\[
P_3=\tfrac14(I+S^3+S^6+S^9),\qquad
P_4=\tfrac13(I+S^4+S^8)
\]
have images of dimensions three and four. If $E=P_3-P_4$, then
\begin{equation}
\frac{\Tr E}{12}=-\frac1{12}.
\label{eq:minus12}
\end{equation}
The scalar comes from a normalized rank difference, not an
ordinary sum of positive integers.

\begin{theorem}[Representation of the channel reader as a ratio of determinants]
Let $\mathcal H_3=\operatorname{Fix}(S^3)$ and
$\mathcal H_4=\operatorname{Fix}(S^4)$. For $0<s<1$,
\begin{equation}
F(s)=\frac{\det_{\mathcal H_3}(I-sS)}{\det_{\mathcal H_4}(I-sS)}
=\frac{1-s^3}{1-s^4}
=\frac{1+s+s^2}{1+s+s^2+s^3}.
\label{eq:det}
\end{equation}
Its continuous extension at $s=1$ equals $3/4$.
\end{theorem}
\begin{proof}
The restriction of $S$ to each image is a cycle of three or four
positions, respectively. Their characteristic polynomials give
$\det(I-sS)=1-s^r$ for $r=3,4$. Canceling the common factor $1-s$
leaves the last expression, regular at $s=1$.
\end{proof}

Substitution $s=q^{30}$ gives
$F(q^{30})=(1-q^{90})/(1-q^{120})$.
The same pair of spaces producing \eqref{eq:minus12}
produces the channel quotient. This equality does not identify
a twelve-coordinate matrix with the entire TPK state.

\subsection{Oriented information and the complex sector}
In the antipodal sector $V_-=\operatorname{Ran}((I-S^6)/2)$,
the operator $J=S^3$ satisfies $J^2=-I$ and $J^{\mathsf T}=-J$.
With $Q=P_4(I-S^6)/2$, one has
\[
\operatorname{rank}Q=2,\qquad E|_{V_-}=-Q,\qquad B|_{V_-}=5I-3Q.
\]
On $\operatorname{Ran}Q$, $S^2=-I$ and the determinant is $1+s^2$.
Thus,
\[
F(s)=\frac{1+s+s^2}{1+s}\,\frac1{1+s^2}.
\]
Orientation reversal changes the sign of the form
$\omega(u,v)=\langle Ju,v\rangle$, but preserves $B$ and $F$.
Scalar evaluation does not distinguish all oriented transports.

The transported trace also produces
\[
a_n^{\rm tr}=\Tr(S^nE)=3\mathbf1_{3\mid n}-4\mathbf1_{4\mid n},
\]
and, for $|s|<1$,
\[
-\log F(s)=\sum_{n\ge1}\frac{a_n^{\rm tr}}n\,s^n.
\]
This sequence of moments links the cyclic decomposition to the
spectral representation of the reader. The arithmetic nature of particular
values, such as Catalan's constant, additionally requires the corresponding
arithmetic analysis.
